\documentclass[11pt,a4paper]{article}

\usepackage[utf8]{inputenc}
\usepackage[T1]{fontenc}
\usepackage[margin=2.5cm]{geometry}
\usepackage{amsmath,amssymb}
\usepackage{booktabs}
\usepackage{hyperref}
\usepackage{enumitem}
\usepackage{xcolor}

\hypersetup{
  colorlinks=true,
  linkcolor=blue!60!black,
  urlcolor=blue!60!black,
  citecolor=blue!60!black
}

\title{\textbf{Monte Carlo Tree Search for Hex and Y}\\
\large Classical UCT--RAVE, Gumbel AlphaZero, and a Unified Y Extension}
\author{Lea Dodeman \and Nestor Antunano Cabrera}
\date{\today}

\begin{document}

\maketitle
\tableofcontents
\newpage

\begin{abstract}
This repository studies Monte Carlo Tree Search across two connection games,
Hex and Y, inside one shared implementation framework. The project has three
active branches: a classical Hex baseline reproducing the UCT+RAVE setting of
\emph{Monte-Carlo Hex}, a learned Hex agent based on Gumbel AlphaZero, and a
new classical implementation for Y. The main engineering goal is consistency:
the classical branches expose the same Python and Cython APIs, use incremental
connectivity tracking, and can be compared without conflating algorithmic
differences with implementation asymmetries. The main research takeaway is
twofold. First, the classical Hex results preserve the qualitative behavior
reported in the literature even with simpler rollouts. Second, the new Y
implementation now matches the Hex branch in logic and speed infrastructure.
Its full-size experiments show the expected simulation-scaling curve, a strong
preference for $C_{\mathrm{uct}}=0$, and a clear benefit from connectivity-aware
rollouts, while also making clear that Y still lacks an analogue of Hex bridge
knowledge with equally strong structural impact.
\end{abstract}

\section{Introduction}

Monte Carlo Tree Search (MCTS) is a search procedure that incrementally builds
a partial game tree through four repeated phases: selection, expansion,
simulation, and backpropagation. Rather than expanding all children uniformly,
it allocates simulations to the lines of play that appear most promising under
the current statistics. This makes MCTS particularly effective in deterministic
perfect-information games with large branching factors and difficult static
evaluation functions.

This repository began as a Hex project and then expanded into a broader study
of connection-game search. It now contains:

\begin{itemize}[leftmargin=1.5em]
    \item \texttt{HexClassic/}: a classical UCT+RAVE implementation inspired by
    Cazenave and Saffidine's \emph{Monte-Carlo Hex} \cite{cazenave2010montecarlohex},
    \item \texttt{HexGumbel/}: a neural-network-guided Hex agent using the
    Gumbel AlphaZero root policy \cite{danihelka2022gumbel},
    \item \texttt{YClassic/}: a classical Y implementation designed to reach the
    same logic and performance standard as \texttt{HexClassic/}.
\end{itemize}

The documentation is organized around the repository as a whole rather than as
three disconnected folders. Shared engineering choices are described once, then
each branch is discussed only where its algorithmic behavior or game-specific
logic diverges.

\section{Games and Scope}

\subsection{Hex}

Hex is played on an $N \times N$ rhombus of hexagonal cells. One player tries
to connect top to bottom, while the other tries to connect left to right.
Players alternately occupy empty cells, and stones are never moved or removed.
Hex has no draws on a full board, which makes it a clean fit for simulation-
based search.

\subsection{Y}

Y is played on a triangular board with hex-style adjacency. A player wins by
forming one connected component that touches all three sides of the triangle.
Corner cells count as belonging to both adjacent sides, and a full board has
exactly one winner \cite{boardspace_y,gambiter_y}. Standard descriptions of Y
also note that the game can be played either on a regular triangular hex grid
or on the ``official'' Schensted/Titus board that inserts three five-neighbor
points \cite{boardspace_y,gambiter_y}. This project uses the regular triangular
grid because it is the simplest variant for a direct implementation comparison
with Hex.

\subsection{Repository Modules}

\begin{center}
\begin{tabular}{@{}lll@{}}
\toprule
Module & Game & Role \\
\midrule
\texttt{HexClassic/} & Hex & Classical UCT+RAVE baseline and reproduction study \\
\texttt{HexGumbel/} & Hex & Learned policy-value search with Gumbel root search \\
\texttt{YClassic/} & Y   & Classical Y search adapted from the Hex baseline \\
\bottomrule
\end{tabular}
\end{center}

The Y work in the current repository is focused on the classical path only.

\section{Search Methods}

\subsection{Classical MCTS with UCT and RAVE}

The classical baseline follows the UCT+RAVE structure of
\cite{cazenave2010montecarlohex}. Selection uses the empirical action value
augmented by the UCT exploration bonus:
\[
\mathrm{UCT}(s,a) = Q(s,a) + C_{\mathrm{uct}}
\sqrt{\frac{\log N(s)}{N(s,a)}}.
\]

RAVE (Rapid Action Value Estimation), also called AMAF, shares information
between branches by treating moves played later in a simulation as if they had
been tried immediately. The implementation interpolates between the Monte Carlo
estimate and the RAVE estimate using the same bias-controlled blending scheme
used in \cite{cazenave2010montecarlohex}. This is important in early search,
where raw visit counts are still sparse.

\subsection{Gumbel AlphaZero}

The learned Hex branch replaces rollout evaluation with a neural policy-value
model. At interior nodes it uses the PUCT rule
\[
Q(s,a) + c_{\mathrm{puct}} P(s,a)\frac{\sqrt{N(s)}}{1+N(s,a)},
\]
and at the root it follows the Gumbel AlphaZero policy-improvement mechanism of
\cite{danihelka2022gumbel}. Legal moves receive Gumbel perturbations, the root
policy is filtered to a candidate set, and sequential halving allocates the
simulation budget among the strongest contenders. Relative to a pure AlphaZero
pipeline \cite{silver2017mastering}, the project adds supervised pretraining
from classical Hex labels because the available compute budget is limited.

\section{Unified Implementation Strategy}

The repository uses one shared engineering pattern across the classical
branches, and this is the main reason the new Y implementation could be brought
up to the same standard as the older Hex code quickly.

\subsection{Reference Path and Fast Path}

Each classical module exposes:

\begin{itemize}[leftmargin=1.5em]
    \item a pure Python reference implementation that is easy to inspect,
    \item a Cython implementation with the same public API for serious
    experimentation.
\end{itemize}

This keeps the algorithm readable while still making large experimental runs
practical.

\subsection{Incremental Connectivity}

Both Hex and Y are connection games, so the critical inner-loop operation is
not piece movement but component maintenance. The repository therefore tracks
connectivity incrementally with Union-Find structures rather than recomputing
graph connectivity after every move.

In Hex, components connect to two virtual borders. In Y, components instead
carry a three-bit side mask. A Y move wins as soon as the merged component mask
becomes \texttt{111}, meaning that the connected component touches left, right,
and bottom simultaneously.

\subsection{Cython Acceleration}

The fast classical backends follow the same optimization pattern once and only
once across the whole project:

\begin{itemize}[leftmargin=1.5em]
    \item board state in typed C arrays,
    \item board cloning with \texttt{memcpy},
    \item pooled MCTS nodes rather than Python object churn,
    \item flat RAVE arrays instead of Python dictionaries in hot loops,
    \item explicit RNG seeding for multiprocessing workers.
\end{itemize}

The point is not that Hex and Y require two different stories about why Cython
helps. They benefit from the same structural reduction in interpreter overhead.

\section{HexClassic}

\subsection{Goal}

\texttt{HexClassic/} is a comparative study rather than an attempt to recreate
the strongest known Hex engine. It asks whether the qualitative findings of
\emph{Monte-Carlo Hex} still hold if the rollout policy is simplified from the
paper's type-3 setting to type-2 bridge defense only.

\subsection{Results}

All results below use the Cython backend on $11 \times 11$ boards with 200
games per data point.

\paragraph{Simulation count.}

\begin{center}
\begin{tabular}{@{}rrrr@{}}
\toprule
Simulations & Ours & Paper & Delta \\
\midrule
1{,}000  & 31.5\% & 6.0\%  & +25.5 \\
2{,}000  & 45.5\% & 11.5\% & +34.0 \\
4{,}000  & 41.5\% & 20.0\% & +21.5 \\
8{,}000  & 50.0\% & 33.0\% & +17.0 \\
32{,}000 & 55.5\% & 61.0\% & -5.5  \\
64{,}000 & 66.5\% & 68.5\% & -2.0  \\
\bottomrule
\end{tabular}
\end{center}

The expected monotone trend is preserved. The weaker rollout policy compresses
the low-simulation part of the curve, but the results converge toward the paper
once the tree policy dominates the playout quality.

\paragraph{UCT exploration constant.}

\begin{center}
\begin{tabular}{@{}rrrr@{}}
\toprule
$C_{\mathrm{uct}}$ & Ours & Paper & Delta \\
\midrule
0.0 & 88.0\% & 61.0\% & +27.0 \\
0.1 & 62.5\% & 60.0\% & +2.5  \\
0.2 & 52.5\% & 55.5\% & -3.0  \\
0.4 & 43.0\% & 42.0\% & +1.0  \\
0.5 & 55.5\% & 41.0\% & +14.5 \\
0.6 & 47.5\% & 35.5\% & +12.0 \\
0.7 & 50.5\% & 32.5\% & +18.0 \\
\bottomrule
\end{tabular}
\end{center}

The strongest setting remains $C_{\mathrm{uct}} = 0$ when RAVE is enabled.
This matches the interpretation from \cite{cazenave2010montecarlohex}: RAVE is
already doing the work of early exploration.

\paragraph{Rollout policy and RAVE bias.}

\begin{center}
\begin{tabular}{@{}lrr@{}}
\toprule
Experiment & Result & Interpretation \\
\midrule
Random vs bridge rollouts & 28.5\% & Bridges clearly stronger than random \\
$b = 0.0005$ vs $0.001$ & 47.0\% & \\
$b = 0.00025$ vs $0.001$ & 49.0\% & RAVE tuning is flatter than in the paper \\
$b = 0.000125$ vs $0.001$ & 46.0\% & \\
\bottomrule
\end{tabular}
\end{center}

The main conclusion is that the UCT+RAVE core remains robust even when the
rollout policy is intentionally simplified.

\section{HexGumbel}

\subsection{Method}

\texttt{HexGumbel/} uses a ResNet that outputs both a policy over legal moves
and a scalar value estimate. Classical rollouts disappear; the search is guided
by the network and by the Gumbel AlphaZero root policy
\cite{danihelka2022gumbel}. The project uses a supervised warm start from
\texttt{HexClassic/} labels before entering self-play. This makes the training
pipeline more stable than a pure cold-start AlphaZero setup under limited
compute.

\subsection{Results}

The learned agent is evaluated on $7 \times 7$ Hex with only 16 simulations per
move against a classical Hex baseline using 16{,}000 simulations per move.

\begin{center}
\begin{tabular}{@{}rrrr@{}}
\toprule
Iteration & vs Random & vs MCTS 16K & Buffer Size \\
\midrule
5  & 100\% & 10.0\% &  69{,}220 \\
10 & 100\% & 32.5\% & 132{,}931 \\
20 & 100\% & 45.0\% & 255{,}492 \\
30 & 100\% & 55.0\% & 121{,}818 \\
40 & 100\% & 62.5\% &  63{,}061 \\
50 & 100\% & 55.0\% &  61{,}491 \\
60 & 100\% & 50.0\% &  62{,}208 \\
70 & 100\% & 70.0\% & 188{,}366 \\
\bottomrule
\end{tabular}
\end{center}

The learned agent eventually reaches roughly 65--70\% against the classical
baseline despite using a thousand times fewer simulations. The result is not
that search became unnecessary; it is that a better prior and a learned value
allow the search budget to be deployed much more selectively.

\section{YClassic}

\subsection{Design Goal}

The Y branch started as a draft and lacked the implementation maturity of
\texttt{HexClassic/}. The goal was therefore not just to make it run, but to
bring it to the same standard in three dimensions:

\begin{itemize}[leftmargin=1.5em]
    \item rule correctness,
    \item architectural parity with the Hex baseline,
    \item practical speed through a Cython backend.
\end{itemize}

\subsection{Rule Model and Board Logic}

The regular-triangle Y implementation uses
\[
n = \frac{s(s+1)}{2}
\]
cells for side length $s$. Each move belongs to zero, one, or two sides of the
triangle depending on whether it lies in the interior, on an edge, or in a
corner. For each player, every connected component stores a three-bit side
mask:
\[
\texttt{LEFT} = 001,\quad
\texttt{RIGHT} = 010,\quad
\texttt{BOTTOM} = 100.
\]

Whenever a new stone is placed, it merges with same-color neighboring
components and the masks are OR-combined. A player wins as soon as one merged
component reaches \texttt{111}. This makes Y win detection constant-time after
each move.

\subsection{Search Adaptation}

Y does not have a direct analogue of Hex bridge templates. The classical Y
rollout policy therefore uses a local connectivity score rather than a fixed
bridge-save response. Candidate rollout moves are preferred when they:

\begin{itemize}[leftmargin=1.5em]
    \item merge multiple friendly components,
    \item increase the number of sides touched by the merged component,
    \item occupy strategically relevant cells near enemy groups.
\end{itemize}

The search core itself remains the same as in Hex: one UCT+RAVE selection loop,
AMAF backpropagation, and full-board playouts.

\subsection{Verification}

The Y implementation was checked at three levels:

\begin{enumerate}[leftmargin=1.5em]
    \item exhaustive full-board checks on small triangular boards to confirm
    that every coloring has exactly one winner,
    \item random-game checks against a naive connectivity evaluator,
    \item Python-vs-Cython consistency checks on identical move sequences.
\end{enumerate}

The automated regression suite currently passes in both Python-only mode and
with the compiled Cython backend.

\subsection{Performance}

An empty-board benchmark on side length 7 with 2{,}000 simulations produced the
following move-selection times:

\begin{center}
\begin{tabular}{@{}lr@{}}
\toprule
Backend & Time \\
\midrule
Python & 0.298 s \\
Cython & 0.0135 s \\
Speedup & 22.1$\times$ \\
\bottomrule
\end{tabular}
\end{center}

This closes the original engineering gap between the Y draft and the mature
Hex implementation. On the larger side-length-11 benchmark with the Cython
backend and the same 2{,}000-simulation root search, Hex took 0.0276 s and Y
took 0.0377 s. In other words, Y is only about 1.37$\times$ slower than Hex in
the matched non-toy setting.

\subsection{Full-Scale Experimental Results}

The full Y experiment suite was then run with the same protocol shape used in
\texttt{HexClassic/}: Cython backend, side length 11, 200 games per point, and
seed 42. The experiment runner saves one JSON artifact per table plus a
consolidated \texttt{all\_results.jsonl} file in \texttt{YClassic/results/}.

\paragraph{Simulation count.}

\begin{center}
\begin{tabular}{@{}rr@{}}
\toprule
Simulations & Win rate against 16{,}000-sim reference \\
\midrule
1{,}000  & 2.5\% \\
2{,}000  & 17.0\% \\
4{,}000  & 34.0\% \\
8{,}000  & 44.5\% \\
32{,}000 & 55.0\% \\
64{,}000 & 60.5\% \\
\bottomrule
\end{tabular}
\end{center}

This is the expected qualitative behavior: more simulations still mean stronger
play, and the curve crosses the 50\% mark once the search budget exceeds the
16{,}000-simulation reference.

\paragraph{UCT exploration constant.}

\begin{center}
\begin{tabular}{@{}rr@{}}
\toprule
$C_{\mathrm{uct}}$ & Win rate against $C_{\mathrm{uct}}=0.3$ \\
\midrule
0.0 & 90.0\% \\
0.1 & 76.0\% \\
0.2 & 59.0\% \\
0.4 & 45.5\% \\
0.5 & 43.5\% \\
0.6 & 45.5\% \\
0.7 & 47.5\% \\
\bottomrule
\end{tabular}
\end{center}

The Y branch shows the same qualitative pattern seen in Hex: when RAVE is
active, the strongest setting is again $C_{\mathrm{uct}}=0$, and the preference
is even more decisive here.

\paragraph{Rollout policy and RAVE bias.}

\begin{center}
\begin{tabular}{@{}lrr@{}}
\toprule
Experiment & Result & Interpretation \\
\midrule
Random vs connectivity rollout & 37.5\% & Connectivity is clearly stronger \\
$b = 0.0005$ vs $0.001$ & 51.5\% & \\
$b = 0.00025$ vs $0.001$ & 50.0\% & RAVE bias is comparatively flat \\
$b = 0.000125$ vs $0.001$ & 47.5\% & \\
\bottomrule
\end{tabular}
\end{center}

The important point is that the Y branch no longer only ``looks plausible'' in
small calibration tests. At the same non-toy 11-side scale used for the Hex
study, it shows clear simulation scaling, a strong tree-policy preference, and
a measurable gain from rollout knowledge. What remains flatter than in Hex is
not the implementation quality, but the strength gap between lightweight Y
heuristics and the stronger bridge knowledge available in Hex.

\section{Discussion}

Taken together, the three modules show three complementary points.

\begin{itemize}[leftmargin=1.5em]
    \item Classical search remains competitive and interpretable in connection
    games, especially when paired with well-chosen rollout knowledge.
    \item Learned search can dominate rollout-based evaluation when the policy
    and value network are good enough, even under a much smaller simulation
    budget.
    \item Extending a mature search system to a new game is not only a matter
    of copying code. The board geometry, win condition, and available domain
    heuristics all matter. Y could inherit the Hex engineering backbone, but it
    still required new logic, and its heuristics remain weaker and less
    structured than Hex bridge knowledge even after the full-scale run.
\end{itemize}

The most important engineering lesson is that the shared architecture paid off.
Once the repository had a consistent pattern for Python references, Cython
fast paths, Union-Find connectivity, and reproducible experiments, the Y branch
could be upgraded quickly without becoming a one-off implementation.

\section{Conclusion}

This repository now contains a coherent MCTS study across Hex and Y. The
classical Hex branch reproduces the main qualitative findings of
\emph{Monte-Carlo Hex}. The learned Hex branch shows that Gumbel AlphaZero can
outperform a far more heavily simulated classical baseline on the chosen 7x7
setting. The Y branch is now no longer a draft: it has correct rule handling, a
matching Python/Cython structure, automated regression tests, and a compiled
backend with a measured 22x speedup over the Python reference.

More importantly, the Y branch is now documented with real non-toy experiments
rather than only pilot checks. On 11-side boards it shows the expected
simulation-scaling curve, strongly prefers $C_{\mathrm{uct}}=0$ under RAVE, and
its connectivity rollout defeats purely random playouts by a clear margin.
Performance is also in the right regime: at matched 11-side, 2{,}000-simulation
root search, Y is only about 1.37$\times$ slower than Hex.

The natural next step for Y is not more implementation cleanup but stronger
domain knowledge: the search stack is ready, and the remaining open question is
which rollout or prior knowledge best captures Y-specific structure.

\bibliographystyle{abbrv}
\bibliography{references}

\end{document}
